
\DocumentMetadata{tagging=on,
	tagging-setup={math/setup=mathml-SE},
	pdfstandard=ua-2,
	lang=en-US
}
\documentclass{article}

%Math Libraries
\usepackage{multirow, longtable, amsmath, amssymb, amsthm}
\usepackage[mathscr]{euscript}

%Formatting
\usepackage[margin = 1 in]{geometry}
\usepackage{indentfirst}
\usepackage{graphicx}
\usepackage{fancyhdr}
\usepackage{tabularx}
\usepackage{cite}

%Diagram Packages
\usepackage{tikz}
\usetikzlibrary{automata,positioning,arrows}

%Standard Commands
\newcommand{\mm}{\mathscr}
\newcommand{\B}{{\mathcal{B}}}
\newcommand{\A}{{\mathcal{A}}}
\newcommand{\N}{{\mathbb{N}}}
\newcommand{\R}{{\mathbb{R}}}
\newcommand{\C}{{\mathbb{C}}}
\newcommand{\Q}{{\mathbb{Q}}}
\newcommand{\Z}{{\mathbb{Z}}}


\newcommand{\abs}[1]{{|{#1}|}}
\newcommand{ \norm }[1]{{ || {#1}||}}
\newcommand{\cl}[1]{{\overline{#1}}}
\newcommand{\pd}{\partial}
\newcommand{\ip}[2]{ \langle {#1},{#2}\rangle }
\newcommand{\ls}{\lesssim}
\newcommand{\gs}{\gtrsim}
\newcommand{\supp}{\text{supp}}

%Result Environment
\newcounter{result}[section]
\newenvironment{res}[1][]{\refstepcounter{result}\par\medskip
	\noindent \textbf{[\thesection.\theresult] #1} \rmfamily}{\medskip}

\newenvironment{defn}[1][]{\refstepcounter{result}\par\medskip \noindent \textbf{Definition [\thesection.\theresult] #1} \rmfamily}{\medskip}

%Proof Environment
\newenvironment{pf}{%
	\par%
	\medskip
	\leftskip=2em\rightskip=2em%
	\noindent\ignorespaces \textit{Proof:} \newline}{%
	\,\, \newline \textit{Q.E.D.}\par\medskip}

\title{Math 126 Summer 2026 HW 2}
\author{Lec 001, Ebbighausen}


\begin{document}
	\pagestyle{fancy}
	\fancyfoot{} 
	\fancyfoot[L]{Math 126}
	
	\maketitle
	
	\textbf{Problem 1:} (Borthwick 4.1, modified) Suppose $u(t,x)$ is a solution to the wave equation 
	\begin{align*}
		\pd_{t}^{2}u - c\pd_{x}^{2}u = 0
	\end{align*}
	
	Part A.) Assume $u(0,x)$, $u(t/2,x+ c\frac{t}{2})$ and $u(t/2, x - c\frac{t}{2})$ are known values. Write a formula for $u(t,x)$ given these three values (Hint: Draw a picture, and consider the domain of dependence of each point). 
	
	Part B.) Assume that $v(t,x)$ satisfies 
	\begin{align*}
		\pd_{t}^{2}v - \pd_{x}^{2}v = 1
	\end{align*}
	and $v(0,2) = 1$, $v(1,3) = v(1,1) = 2$. What is $v(2,2)$?
	
	\textbf{Problem 2:} Show that if $u(t,x)$ satisfies 
	\begin{align*}
		\pd_{t}^{2}u - 2\pd_{t}u - 2 \pd_{x} u - \pd_{x}^{2} u =0
	\end{align*}
	then $w(t,x) = e^{x-t}u$ satisfies the wave equation. Find a formula for $u$ give initial conditions $(u,u_{t})|_{t=0} = (g,h)$. 
	
	
	\textbf{Problem 3:} (Borthwick 4.7) Let $u(t,x)$ solve the Schr{\"o}dinger equation 
	\begin{align*}
		\begin{cases}
			\pd_{t} u - i \Delta u = 0 \\
			u(0,x) = g(x)
		\end{cases}
	\end{align*}
	for $g(x) \in C_{c}^{\infty}(\R)$. 
	
	Part A.) Show that $\int_{\R^{n}} |u|^{2} dx = \int_{\R^{n}} |g|^{2} dx$ (assume that both integrals are defined, and that the Leibniz rule and Green's identities are justified).
	
	Part B.) Show the solution $u$ is uniquely determined by the initial condition $g$. 
	
	
	\textbf{Problem 4:}
	
	Consider the equation 
	\begin{align*}
		t \pd_{t} u - (x+1) \pd_{x} u = 0
	\end{align*}
	
	Part A.) Use the method of separation of variables to form a family of solutions.
	
	Part B.) Given initial condition $u(1,x) = x$, solve the equation using the method of characteristics. Are the two solutions the same?
	
	
	Part B.) 
\end{document}